% ============================================================================
% numbers_perclass.tex — added 2026-08-18 for main_revised.tex Sec. 2.5
%
% Per-class decomposition of the softmax-vs-prototype contrast, supporting
% Supplementary Figs. S1 and S2. Computed by
%   scripts/paper_figs/figS1_perclass_head.py
%   scripts/paper_figs/figS2_redistribution.py
% from results/paper_reruns/predictions/{B1,B2}_corrected__pointnet2__
% {softmax,proto}__subject__seed{0..4}.npz — five-seed mean per class.
%
% VALIDATION: recomputing from the raw prediction files reproduces every
% published macro value exactly, which is why these figures can be trusted to
% decompose them:
%   33-class  F1 0.280 -> 0.341 (\dxFoneSubjSM/\dxFoneSubjPr),
%             AUC 0.913 -> 0.897 (\dxAucSubjSM/\dxAucSubjPr)
%   19-class  F1 0.326 -> 0.349, AUC 0.811 -> 0.812
% ============================================================================

% ── The "unchanged" band. READ THIS BEFORE QUOTING 14/9/10. ─────────────────
% The manuscript's counts of improved / reduced / unchanged classes are NOT a
% comparison against exact zero. Exact comparison gives 15 up, 11 down, 7 flat.
% The published 14/9/10 corresponds to treating |dF1| <= 0.01 as unchanged, and
% that threshold was previously undeclared anywhere in the text. It is now
% stated in the prose and drawn on Fig. S1(a). Do not change the counts without
% changing the band, and do not quote them without the band.
\newcommand{\headBand}{0.01}
\newcommand{\headClassesUpB}{14}
\newcommand{\headClassesDownB}{9}
\newcommand{\headClassesFlatB}{10}
% Of the 10 unchanged classes, 7 score F1 = 0 under BOTH decision rules: the
% model never places them first either way, so they are floor-bound rather than
% genuinely tied. This distinction is the substantive part of the "no change"
% count and is drawn as a separate series in Fig. S1(a).
\newcommand{\headClassesZeroBoth}{7}
\newcommand{\headClassesTrueTie}{3}

% ── Magnitude of redistribution vs change in discrimination ────────────────
% Mean absolute per-class change, prototype minus softmax. The ratio is the
% quantitative form of "changed which diagnoses were classified well, while
% leaving overall discrimination similar".
\newcommand{\dxPerClassFoneMad}{0.091}
\newcommand{\dxPerClassAucMad}{0.038}
\newcommand{\dxPerClassRatio}{2.4}
\newcommand{\synPerClassFoneMad}{0.045}
\newcommand{\synPerClassAucMad}{0.026}
\newcommand{\synPerClassRatio}{1.8}

% ── Macro AUC under the head swap, prototype minus softmax ─────────────────
% !! SIGN. contrast_tests.csv reads SOFTMAX MINUS PROTOTYPE; these are negated
% to match the prose convention, exactly as \dxHeadAccDiff etc. in numbers.tex.
% contrast_tests.csv | C1a_head_subject[pointnet2]
%
% The two tasks do NOT license the same statement and must not be merged:
%   19-class  emitter verdict "TIE CONFIRMED, tight (CI within +/-2%)"
%   33-class  emitter verdict "TIE NOT ESTABLISHED (CI contains 0 but is too
%             wide to exclude a material difference)"
% So discrimination is demonstrably preserved at 19 classes and merely not
% distinguishable at 33. Writing "unchanged" for both would overclaim the 33.
% \synHeadAucDiff, \synHeadAucCI, \dxHeadAucDiff and \dxHeadAucCI ALREADY EXIST
% in numbers.tex (lines ~272 and ~286) and are NOT redefined here. Recomputing
% them independently from the raw prediction files gave +0.1 / -1.6 pp with the
% same intervals, i.e. the existing macros check out. numbers.tex lines 280-285
% already carry the warning that only the 19-class AUC null is tight; the prose
% added for Sec. 2.5 follows that rule rather than restating it.
%
% Only the two verdict phrases are new, so that the distinction is carried by a
% macro instead of being re-worded each time it is needed.
\newcommand{\synHeadAucVerdict}{confirmed as equivalent within a $\pm2$ percentage-point margin}
\newcommand{\dxHeadAucVerdict}{not distinguishable from zero, though the interval is too wide to establish equivalence}
